% !TeX root = ../papers/nc-p3.tex
\section{六次切断の線形行列と特異三次核曲面}
\label{sec:cofactor-eighteen-cubic-incidence}
\begin{theorem}\label{thm:cofactor-eighteen-normal-cubic}
$d=18$ では正規三次曲面は核像にならない。非正規三次曲面でも、
一般核切断が算術種数四である場合、または種数三で入力逆写像が特異直線の一般点において
三次曲面上の正則写像となる場合は生じない。
\end{theorem}
\begin{proof}[算術種数四の切断]
核写像は三次曲面 $K$ に双有理である。$J$ が種数四なら $h^0(\mathcal O_J(2))=9$ なので
二次曲面に含まれ、次数比較で $J$ は $K$ との完全交叉である（共通成分なしの純な完全交叉で次数六を使い切る）。
その二次式は一意である。実際二つの異なる二次曲面の共通曲線は次数四以下、また $J$ は整非平面である。

逆写像は $K$ 上の四つの二次式で与えられる。非正規な場合も含めこれを示す。
一般入力平面の係数 $t$ を変数にすると一意な二次式 $Q_t$ の係数は有理的に定まる
（関数体上の一般切断の二次イデアルの一次元空間）。斉次化して共通因子を除く。
$K\times\mathbb P^{3*}$ 上の $Q_t(y)=0$ には一般 $t$ 上で既約な被約切断 $J_t$ が一重に入る。
ほかの因子成分は $t$ の超曲面の上の積全体でなければならず、$K$ 上で恒等零になる
非零二次式はないので共通因子の除去に反する。一般の $y$ 上でこの入射因子は
入力逆像を通る平面の一次式であり、その一般点では一重である。従って $t$ の次数は一となり、
四係数が所要の入力二次式である。共通曲線があれば一般 $J_t$ の成分になり不可能である。

$G\to J$ の唯一の余長一点 $q$ は逆基点である。一般平面ではその上の各入力点は階数三であり、
$G$ はそこで平面曲線に同型である。補題\ref{lem:cofactor-colength-one-multiplicity}により
$J$ の $q$ での重複度は二となる。$K$ の重複度は高々二である。
$S_1=\operatorname{Bl}_q K$ を用いる。$K$ が $q$ で滑らかなら一般逆切断は位数二で、
固有変換 $J'$ は $A'=2L-2e$ のCartier因子である（$e$ は例外因子）。
$K$ の重複度が二なら一般逆切断は位数一で、その先頭一次式は $K$ の接錐のどの成分も含まない。
実際一般平面での局所交点数が二であり、位数積の下界は二、接方向を共有すれば不等号が厳しくなる。
従ってこのとき $J'$ は $A'=2L-e$ のCartier因子である。
$S_1$ は空間の一点ブローアップ内の固有変換超曲面であり局所完全交叉である。
両場合とも $J'\to J$ は有限で、$J'$ の算術種数は三である。
滑らかな点なら二重基点の随伴公式、二重点なら $K_{S_1}=-L$, $L^2=3$, $Le=0$, $e^2=-2$ を使う。
これらの公式は超曲面の随伴で非正規でも成立する。よって $J'=G$ である。

滑らかな基点なら $A'^2=8$ で残余基点なし、入力射は $2(L-e)$ で与えられる。
正次元繊維は $q$ を通る直線だけで、高々三本が接平面に含まれる。
一般の滑らかな三次曲面とその一点への変形で入力束は平面七点の $6h-2\sum e_i$。
二重点なら $h=L-e$ は平面への双有理射、$A'=L+h$ である。
入力比の正則性により残余基点は長さ $A'^2-8=2$ の局所完全交叉である。
$h$ が潰す曲線は $q$ を通る直線の固有変換であり、相異なるものは射影方向が異なり互いに素である。
各々で $A'$ 次数は一なので残余二点に会う高々二本だけが入力にも潰れる。
$h$ で正次数の入力縮約曲線は高々二本である。実際 $A'-h=L$ の切断は四独立で、
長さ二の残余条件後も二独立であり、平面像は次数三である。入力次数零の各曲線は
これらの固定成分に入る。残余ブローアップの例外直線は入力次数一である。

二重点の場合、$S_1$ を空間の一点ブローアップの $3L-2e_{\rm amb}$ の系で動かすと一般繊維は滑らかである。
一般には二重点つき三次の式は $zQ_2+Q_3=0$、平面への逆射影の六基点は
$Q_2=Q_3=0$ の相異なる六点で、$S_1$ はそのブローアップである。
この滑らかなモデルでさらに六中心を一般の平面六点に動かせる。
$h,L$ は平面超平面および六点の反標準類として動き、残余二点も一般にできる。
滑らかな基点の場合は空間のブローアップの $3L-e_{\rm amb}$ の系で平滑化できる。
いずれも $H^1(\mathcal O_{S_1})=H^2(\mathcal O_{S_1})=0$ は空間ブローアップの超曲面列から従う
（随伴で $K_{S_1}=-L$ または $-L+e$、超曲面列の双対も用いる）。
一般切断を除くと $J'=G$ 上の次数八の線束になるので入力四切断も持ち上がる。
局所完全交叉残余基点の持ち上げと三点繊維の持ち上げは定理\ref{thm:cofactor-eighteen-cone-not-birational}と同じである。
一般の平面七点モデルまたは八点モデルで $\chi(L+tH)=4(t+1)^2$。
従って中心の有限入力開集合で $M=f_*L$、少なくとも十七または十六個の孤立三点繊維が得られる。
一般の長さ十四または十三に反する。
\end{proof}

\begin{proof}[算術種数三の切断]
$K$ が正規、または入力逆写像が余次元一で定義されると仮定する。このとき
一般 $J$ は局所平面曲線である。$K$ の滑らかな開集合の基点外ではBertiniを使う。
基点および $K$ の孤立特異点の上では一般平面での入力点が階数三となる。
非正規三次の特異曲線は一本の直線であり、その上の核グラフ解消上の全曲線のうち
入力でも曲線に写るものの像は次数六以下である。実際その直線を含む核超平面の引戻しに入り、
$HL=6$ である。従って一般平面切断はそこで整な十八次の階数二曲線を避ける。
これらの点で $J=G$ は階数三の入力平面曲線に局所同型である。

$J$ は二次曲面に入らない（入れば種数四の完全交叉）。切断次数六が標準次数四を超えるので
$h^0(\mathcal O_J(1))=4$, $h^0(\mathcal O_J(2))=10$。
従って次数一、二の制限は同型、イデアルは三正則であり、全次数で中間コホモロジーが消える。
Hilbert分子は $1-4t^3+3t^4$ である。よってHilbert--Burch行列は一次式の四行三列で、
その極大小行列式が三次生成元である。一つを $K$ の式として選び、対応する三行を $A$、残りの一行を $b$ とする。
全行列の階数は $J$ 上で常に二である。階数一以下なら全生成元が極大イデアルの二乗に入り、
$J$ の埋込み次元が三になってしまう。

$S_1=\{(y,z):A(y)z=0\}\subset\mathbb P^3\times\mathbb P^2$ とし両超平面を $L,h$ とする。
$A$ の階数一以下の点は全て $J$ 上で $K$ の特異点である。有限個しかない
（$J$ は非平面六次、三次の特異軌跡は高々直線と有限点）。階数零は全行列の階数から不可能。
従って三つの $(1,1)$ 式は期待次元の完全交叉を定め、$K$ への射 $p$ は
有限個の例外直線の外で同型である。一般被約かつ完全交叉で余分な曲面成分はなく $S_1$ は整である。
$b(y)z=0$ が切る因子は $J$ と同型である。全行列の唯一の核が小行列式で各点で正則に復元されるからである。

入力切断は $S_1$ 上の $A'=L+h$ の切断に持ち上がる。詳細には、三式のKoszul列に $h$ を足して
第一射影で押すと $p_*\mathcal O_{S_1}(h)$ は三行の一次式の正方表示を持つ
（$\mathbb P^2$ の $\mathcal O(-1),\mathcal O(-2)$ は全コホモロジーが消える）。
よってこの直像および $p_*\mathcal O_{S_1}(L+h)$ は $K$ 上の最大Cohen--Macaulay層である。
$K$ の有限集合を除けば $p$ は同型で、入力写像は仮定によりそこで正則である。
選んだ切断 $J=(bz=0)$ を通して両線束がその開集合で同一視される
（一般入力超平面の引戻しはそこで被約な $J$。特異直線は成分に入らない）。
従って入力四切断は有限集合を越えて直像の切断に延長する。選んだ切断は例外直線を含まないので
持ち上げた系にも固定成分がない。一般のその切断も各例外直線を長さ一で切り、
その像 $J$ への射は同型である。実際各例外直線では新切断も $z$ の一次式であり、
局所的に持ち上げた一行を $A$ に加えた行列の核が一意となる
（$p_*\mathcal O(h)$ の上の三生成表示を使う）。
従って再び残余基点は二生成で、$L^2=3$, $Lh=3$, $h^2=1$ から長さ二である。

$h$ による平面への射は双有理で、繊維は線形である。実際三式を $y$ の一次式として読めばよい。
平面繊維は整三次 $K$ に平面を入れるため不可能である。正次元繊維は相異なる平面点の上の直線で、
$L$ 次数一ゆえ残余二点を通る高々二本だけが入力にも潰れる。
他の入力縮約曲線は前半と同じ平面三次の固定成分の論法で高々二本である。
三つの $(1,1)$ 式を一般に動かせば滑らかな平面六点のブローアップになる
（一般の三行四列の平面一次行列の階数二軌跡は六個の単純点で、核の入射がそのブローアップ）。
これは六点のHilbert--Burch表示により一般の六点配置を含む。
$H^1(\mathcal O_{S_1})=H^2(\mathcal O_{S_1})=0$ はKoszul列から従い、切断と長さ二の基点を再び変形できる。
平面八点の $4h-\sum e_i$ の三点数十三と中心の十六以上が矛盾する。
\end{proof}
\path{computations/verify_cofactor_cubic_incidence.py} は完全交叉の二偏極交点数、Hilbert分子と
特異基点の二つの数値型を検算する。
